% Part: normal-modal-logic
% Chapter: frame-correspondence
% Section: introduction

\documentclass[../../../include/open-logic-section]{subfiles}

\begin{document}

\olfileid{nml}{frd}{int}

\olsection{ভূমিকা}

অভিগম্যতা সম্পর্ক এবং সেই সম্পর্কবিশিষ্ট মডেলে নির্দিষ্ট !!{formula}s-এর সত্যতার মধ্যে কী সম্পর্ক, তা মোডাল যুক্তিবিদদের একটি গুরুত্বপূর্ণ প্রশ্ন। যেমন, ধরা যাক অভিগম্যতা সম্পর্কটি প্রতিবিম্বী, অর্থাৎ প্রত্যেক $w
\in W$-এর জন্য $Rww$। অন্য কথায়, প্রত্যেক জগৎ নিজের কাছ থেকেই অভিগম্য। ফলে কোনো জগৎ~$w$-তে $\Box !A$ সত্য হলে, যেসব অভিগম্য জগতে~$!A$ সত্য হতে হবে, তার মধ্যে $w$
নিজেও পড়ে। অতএব, $\mModel{M}$-এর অভিগম্যতা সম্পর্ক~$R$ প্রতিবিম্বী হলে, যেকোনো জগৎ $w$ এবং যেকোনো সূত্র $!A$-এর জন্য সেখানে $\Box !A
\lif !A$ সত্য হবে (অর্থাৎ, $\Box p \lif
p$ স্কিমা এবং তার সব প্রতিস্থাপন-দৃষ্টান্ত~$\mModel{M}$-এ সত্য)।

কিন্তু বিপরীতটি সত্য নয়। যেমন, $\mModel{M}$-এ $\Box p \lif p$ সত্য হলেই যে $R$ প্রতিবিম্বী হবে, তা নয়। সহজেই এমন একটি অপ্রতিবিম্বী মডেল~$\mModel{M}$ পাওয়া যায়, যার সব জগতে $\Box p \lif
p$ সত্য: একটি মাত্র জগৎ~$w$-বিশিষ্ট মডেল নিই, যেখানে সেটি নিজের কাছ থেকে অভিগম্য নয়, কিন্তু $w \in V(p)$। $p$-এর সত্যমান উপযুক্তভাবে বেছে নিলে, অপ্রতিবিম্বী মডেলেও $\Box !A \lif !A$ সত্য করা যায়।

সমাধান হলো সমীকরণ থেকে চলকের মানায়ন~$V$ বাদ দেওয়া। $V(p)$-তে কোন জগৎগুলি আছে তা নির্বিশেষে, যদি~$\mModel{M}$-এর সব জগতে $\Box p \lif p$ সত্য হওয়া চাই, তবে $R$ প্রতিবিম্বী হওয়া আবশ্যক। কারণ, যেকোনো অপ্রতিবিম্বী মডেলে অন্তত একটি জগৎ~$w$ থাকবে, যার জন্য $Rww$ নয়। যদি $V(p) =
W \setminus \{w\}$ ধরি, তবে $w$ ছাড়া সব জগতে $p$ সত্য; সুতরাং $w$ থেকে অভিগম্য সব জগতেও তা সত্য ($w$ যে $w$-এর কাছ থেকে অভিগম্য নয়, তা নিশ্চিত; আর শুধু $w$-তেই~$p$ মিথ্যা)। অন্যদিকে, $w$-তে $p$ মিথ্যা। তাই $w$-তে $\Box
p \lif p$ মিথ্যা।

এ থেকে মানায়নহীন মডেল-গঠনের জন্য একটি স্বতন্ত্র সংকেত চালু করার প্রয়োজন বোঝা যায়: এদের আমরা \emph{কাঠামো} বলি। কাঠামো $\mModel{F}$ হলো জগৎসমষ্টি ও একটি অভিগম্যতা সম্পর্ক নিয়ে গঠিত যুগল $\tuple{W, R}$। তখন প্রত্যেক মডেল $\tuple{W, R, V}$-কে কাঠামো $\tuple{W, R}$-এর \emph{উপর প্রতিষ্ঠিত} বলা যায়। বিপরীতভাবে, একটি কাঠামো তার উপর প্রতিষ্ঠিত মডেলগুলির শ্রেণি নির্ধারণ করে; আর কাঠামোর একটি শ্রেণি সেই শ্রেণির কোনো-না-কোনো কাঠামোর উপর প্রতিষ্ঠিত মডেলগুলির শ্রেণি নির্ধারণ করে। কোনো কাঠামোতে !!a{formula} \emph{বৈধ} হওয়ার ধারণা $\mModel{F} \Entails !A$ এভাবে সংজ্ঞায়িত করি: $\mModel{F}$-এর উপর প্রতিষ্ঠিত সব $\mModel{M}$-এর জন্য $\mSat{M}{!A}$।

এই সংকেত ব্যবহার করে !!{formula}s এবং কাঠামোর শ্রেণির মধ্যে অনুরূপতা স্থাপন করা যায়। যেমন, $\mModel{F} \Entails \Box p \lif p$
তখনই এবং কেবল তখনই, যখন $\mModel{F}$ প্রতিবিম্বী।

\end{document}
